% Inserir imediatamente após SECAO_PILOTO_GOES_C13_CONDICIONAL_LUCAS.tex.
% Requer: \usepackage{booktabs}

\section{Validação técnica do alinhamento multimodal com C13}
\label{sec:smoke-goes16-c13}

Após o piloto condicional apresentado na Seção~\ref{sec:piloto-goes-c13-condicional}, foi realizado um \emph{smoke test} do pipeline de integração entre o canal C13 do GOES-16, o radar do Sumaré e as observações pluviométricas do AlertaRio. Diferentemente do piloto anterior, esta etapa não seleciona eventos pelo target nem estima desempenho preditivo. Seu objetivo é verificar, em um trecho temporal pequeno e rastreável, que o tensor de satélite pode ser construído, armazenado e consumido pelo mesmo loader usado no treinamento multimodal futuro.

\subsection{Dados, escalas e contrato espacial}
\label{subsec:smoke-c13-escalas}

Foi usado o produto \texttt{ABI-L2-CMIPF} Full Disk do GOES-16, disponibilizado publicamente pela NOAA, com o canal infravermelho C13. O canal é centrado aproximadamente em 10,3~$\mu$m e fornece temperatura de brilho em kelvin. Sua resolução espacial nominal é de aproximadamente 2~km no nadir do satélite. Como o Rio de Janeiro está fora do nadir geoestacionário, a resolução efetiva na superfície é um pouco mais grosseira e não deve ser interpretada como uma malha regular de 2~km sobre o município.

Cada cena C13 foi reprojetada por vizinho mais próximo para a mesma grade azimutal equidistante centrada no radar do Sumaré usada pelo dataset auditável de radar: 128~$\times$~128 pixels. Nessa grade, cada célula possui lado aproximado de 2,17~km. A reprojeção torna C13, radar e targets esparsos das estações compatíveis no mesmo sistema de coordenadas, mas não aumenta a resolução espacial efetiva do satélite.

O tensor produzido possui shape $(N, 128, 128, 1)$, tipo \texttt{float32} e um canal, C13. Os valores são armazenados em kelvin e normalizados pelo divisor fixo 400 no loader de treinamento. Essa normalização não usa estatísticas calculadas nos conjuntos de validação ou teste, evitando vazamento de informação entre splits.

\subsection{Alinhamento temporal e causalidade}
\label{subsec:smoke-c13-tempo}

O produto ABI Full Disk possui cadência nominal de 10~minutos no período avaliado, enquanto o radar foi agregado em janelas de 15~minutos a partir de imagens aproximadamente biminutais. Para cada início de janela do radar, o pipeline seleciona a cena C13 mais recente cujo timestamp de \emph{fim} seja menor ou igual ao início da janela do radar. Cenas com idade, medida a partir de seu fim, superior a 20~minutos são marcadas como indisponíveis. A disponibilidade é registrada separadamente, para que o loader descarte qualquer sequência de entrada que contenha C13 ausente.

O teste foi executado para 1\textsuperscript{o} de janeiro de 2022. Foram geradas 96 janelas de radar de 15~minutos, e todas as 96 receberam uma cena C13 causal segundo essa regra. A diferença entre o início da janela de radar e o fim da cena C13 variou de 0,13 a 5,13~minutos, com média de 2,63~minutos. Como as grades temporais de 10 e 15~minutos não coincidem, os timestamps das cenas selecionadas alternaram intervalos de 10 e 20~minutos. Esse efeito é esperado e está documentado no manifesto de origem de cada frame.

O contrato de causalidade atual deve ser interpretado com precisão. Como uma varredura Full Disk leva alguns minutos, uma cena cujo início antecede o radar pode conter pixels adquiridos depois dele. Por essa razão, a seleção exige que o \emph{fim} da cena seja anterior ou igual ao início da janela de radar. Os timestamps de início e fim são preservados no manifesto para rastreabilidade. Essa regra garante causalidade no nível da cena; uma extensão futura poderá incorporar tempos de aquisição por pixel.

\subsection{Resultados da validação}
\label{subsec:smoke-c13-resultados}

A Tabela~\ref{tab:smoke-c13-2022} resume os resultados. Não houve falhas de leitura dos arquivos NetCDF, nem violações da regra causal implementada. Os arquivos Full Disk foram baixados apenas para um diretório temporário, reprojetados e removidos; o artefato permanente contém somente o memmap C13, a máscara de disponibilidade, timestamps de origem e um manifesto com as chaves dos arquivos da NOAA.

\begin{table}[htbp]
\centering
\caption{Resultado do \emph{smoke test} de C13 em 1\textsuperscript{o} de janeiro de 2022.}
\label{tab:smoke-c13-2022}
\begin{tabular}{lr}
\toprule
\textbf{Indicador} & \textbf{Valor} \\
\midrule
Janelas de radar no escopo & 96 \\
Frames C13 disponíveis no escopo & 96 \\
Cobertura no escopo & 100,0\% \\
Falhas de leitura & 0 \\
Violações da regra causal implementada & 0 \\
Idade C13 mínima / máxima & 0,13 / 5,13 min \\
Idade C13 média & 2,63 min \\
Shape do tensor anual & $(27\,147, 128, 128, 1)$ \\
\bottomrule
\end{tabular}
\end{table}

Embora o tensor anual tenha 27\,147 posições em 2022, somente as 96 posições do dia selecionado foram preenchidas neste teste. Assim, a cobertura global anual de 0,35\% é esperada e não representa falha de aquisição. A auditoria registra separadamente o intervalo processado, permitindo avaliar a cobertura de 100\% dentro do escopo solicitado.

Por fim, o loader multimodal foi executado com cinco instantes de entrada, crop espacial das estações e precipitação lagada. Foram obtidas 92 sequências válidas. Cada exemplo teve shape $(6, 5, 53, 67)$, com seis canais: três canais RGB do radar, um canal C13 e dois canais derivados das observações passadas das estações. Esse resultado confirma a compatibilidade de formato entre as três fontes, mas não mede ganho de previsão. A geração histórica completa e a avaliação cronológica independente continuam necessárias antes de comparar um modelo multimodal com C3.
